研究多项式环时, 利用它的递归性质往往很有用, 因为这样便可对变量的个数使用归纳法. 下一个结果相当显然; 但由于经常要用到它, 我们仍仔细陈述如下.

\begin{Proposition}{多项式环的递归性质}\label{prop:polynomial-ring-recursion}
  有一个典范的环同构
  \[
    R[x_{1},\ldots,x_{n}]=R[x_{1},\ldots,x_{n-1}][x_{n}],
  \]
  其中右侧是系数属于环 $R[x_{1},\ldots,x_{n-1}]$ 的单变量 $x_{n}$ 多项式所成的环.
\end{Proposition}

下面的例子说明 Proposition~\ref{prop:polynomial-ring-recursion} 的主要思想.

\begin{Example}{{\textmd{$\mathbb{Z}[x,y]$} 与 \textmd{$\mathbb{Z}[x][y]$}}}\label{ex:iterated-polynomial-ring}
  考虑多项式
  \[
    f=x^{3}y^{2}+xy^{2}-2xy-x+1\in\mathbb{Z}[x,y].
  \]
  把 $y$ 的指数相同的所有项归在一起, 就可将 $f$ 看作 $\mathbb{Z}[x][y]$ 的元素. 这样写得
  \[
    f=(x^{3}+x)y^{2}+(-2x)y+(-x+1)\in\mathbb{Z}[x][y].
  \]
  作为 $y$ 的多项式, $f$ 的系数是 $x^{3}+x$, $-2x$ 和 $-x+1$, 它们都是 $\mathbb{Z}[x]$ 的元素.
\end{Example}

\begin{BookFlowNote}{%
  两个环同构时, 往往有许多可能的同构. \textquotedbl{}典范\textquotedbl{} 一词表示, 在所有可能的同构中有一个自然的选择.
  记号 $\cong$ 用来表示同构, 而 $=$ 用来表示典范同构.
}
\begin{Proof}\unskip\textbf{ of Proposition~\ref{prop:polynomial-ring-recursion}}\par\noindent
  为证明该命题, 我们来描述这个典范的环同构. 如 Example~\ref{ex:iterated-polynomial-ring} 所示, 它只是把 $R[x_{1},\ldots,x_{n}]$ 中某一个多项式内 $x_{n}$ 的指数相同的项归在一起.
  为了更精确地说明这一点, 考虑下列函数:
  \[
    \begin{aligned}
      \phi:R[x_{1},\ldots,x_{n}]&\to R[x_{1},\ldots,x_{n-1}][x_{n}]\\
      \sum_{\alpha\in\mathbb{N}^{n}}a_{\alpha}x_{1}^{\alpha_{1}}\cdots x_{n}^{\alpha_{n}}
      &\mapsto
      \sum_{d\ge0}\left(\sum_{\substack{\alpha\in\mathbb{N}^{n}\\\alpha_{n}=d}}
        a_{\alpha}x_{1}^{\alpha_{1}}\cdots x_{n-1}^{\alpha_{n-1}}\right)x_{n}^{d}.
    \end{aligned}
  \]
  验证 $\phi$ 是环同构, 见 Exercise~\ref{ex:recursive-ring-isomorphism}.
\end{Proof}
\end{BookFlowNote}

下面关于整环上的多项式环的结果, 是 Proposition~\ref{prop:polynomial-ring-recursion} 的第一个应用.
